\bookchapter{第三十四章：車機互聯深度指南——手機與汽車嘅融合}

\chapterintro{汽車正由「四個轆加梳化」進化為「四個轆加智能手機」。無論你揸嘅係燃油車定新能源，手機同車機嘅關係都決定咗你嘅駕駛體驗。呢一章帶你釐清車機互聯嘅三大路線、各派系車機系統嘅特點，以及嗰啲等揸車更順手嘅小技巧。}

\booksection{一、車機互聯嘅三大路線}

當前手機與汽車「聯網」主要有三種思路，理解咗佢哋就理解咗整個生態：

\begin{itemize}
  \item \textbf{手機投屏到車機}：車機淨係當一塊「顯示器」同「音箱」，運算同地圖數據都來自手機。代表方案：Apple CarPlay、Google Android Auto、華為HiCar、百度CarLife。優點係手機生態直接搬上車，唔使重新裝App；缺點係車機本身嘅硬件能力被閒置。
  \item \textbf{手機與車機深度融合}：車機係獨立系統，但同手機共享帳號、應用、數據，甚至能「接力」任務。代表方案：鴻蒙座艙（問界/智界）、比亞迪DiLink、蔚來NIO OS、理想車機系統。優點係體驗一體化；缺點係通常綁定同一品牌生態。
  \item \textbf{CarPlay 2.0深度集成版}：2025年蘋果發布嘅新一代CarPlay，唔再淨係「投屏」，而係接管車機嘅儀錶板、空調控制、座椅調節等核心功能。法拉利、阿斯頓·馬田、保時捷等品牌率先搭載，標誌車機系統進入「手機廠商主導」嘅新階段。
\end{itemize}

\booksection{二、CarPlay / Android Auto / HiCar 對比}

三大主流投屏方案各有側重，選車或選手機時可以對照下表：

\begin{center}
\begin{tabularx}{\linewidth}{|X|X|X|X|X|}
\hline
\textbf{方案} & \textbf{支援手機系統} & \textbf{支援車型} & \textbf{功能豐富度} & \textbf{國內可用性} \tabularnewline
\hline
Apple CarPlay & 僅 iPhone & 極廣，絕大多數車型支援 & 高，地圖/音樂/電話/消息均可 & 好，原生可用 \tabularnewline
Android Auto & Android（國內手機需裝服務框架） & 較廣，合資/進口車居多 & 高，與CarPlay相當 & 一般，需科學上網激活 \tabularnewline
華為HiCar & 華為手機為主，部分支援其他安卓 & 國產車與部分合資車型 & 較高，可與鴻蒙座艙聯動 & 好好，本土生態 \tabularnewline
百度CarLife & iOS 與 Android 均可 & 較廣，舊車型後裝常見 & 中等，更新較慢 & 好好，本土生態 \tabularnewline
\hline
\end{tabularx}
\end{center}

\booksection{三、鴻蒙座艙：手機與車機嘅無縫流轉}

如果講CarPlay係「機投到車機」，鴻蒙座艙則係「手機同車機本來就係一家」。華為手機與問界/智界車機嘅協同，代表咗當前國內車機互聯嘅最高水準：

\begin{itemize}
  \item \textbf{超級桌面}：車機直接調用手機上嘅App，唔使喺車機重新安裝。手機入面便喺聽嘅音樂、喺導航嘅路線，上車自動續接到車機螢幕，落車再返到手機，全程無縫。
  \item \textbf{應用接力}：手機上設好嘅導航，埋車門時車機自動接管；到達停車場，落車後手機繼續步行導航到最終目的地。「最後一公里」唔再斷檔。
  \item \textbf{手機車機共用帳號}：華為帳號一登，車機自動同步你嘅音樂會員、地圖收藏、日程、聯絡人。車機同手機係同一身份嘅唔同終端，而唔係兩套獨立系統。
\end{itemize}

\booksection{四、造車新勢力嘅車機系統}

國產新能源車嘅車機已經由「導航+音樂」進化為移動智能終端，各家嘅玩法各有特色：

\begin{itemize}
  \item \textbf{蔚來NIO OS}：NOMI語音助手係行業標杆，能識別多指令、上文下連續對話，仲帶情感表情。NIO Phone則將手機做成「車匙+第二屏」，深度綁定蔚來生態。
  \item \textbf{理想車機}：多屏交互係亮點，副駕屏可獨立睇劇、仔屏可睇卡通。「任務大師」等你自訂場景，比如「開門自動暫停音樂、埋門繼續播放」。
  \item \textbf{小米澎湃車機}：同小米手機、手錶、家電深度聯動。手機上設嘅鬧鐘會喺車機提醒；車上導航返屋企，屋企冷氣提前開啟。HyperOS打通咗「人車家全生態」。
  \item \textbf{比亞迪DiLink}：標誌性嘅旋轉屏，豎屏睇短影片、橫屏睇地圖，自適應App方向。自帶應用市場，生態豐富度喺傳統車企中領先。
\end{itemize}

\booksection{五、手機當車機——用得嗎？}

如果你嘅車機太老或太難用，能否直接用手機「頂上」？答案係：可以，但要睇場景。

\begin{itemize}
  \item \textbf{摩托羅拉Ready For / 桌面模式}：部分摩托羅拉、三星、華為手機支援「桌面模式」，連上顯示器就係一臺簡易電腦。理論上接車機屏都能用，但實際相容性參差不齊，且缺少汽車專屬交互（如方向盤按鍵）。
  \item \textbf{手機支架 vs 原車車機}：如果原車車機卡頓、地圖唔更新，一個穩固嘅手機支架+手機導航係務實之選。但要注意：(1)手機螢幕細，揸車時睇唔清；(2)夏天暴曬會等手機過熱降頻甚至自動關機；(3)長時間導航耗電快，需配車載充電器。
  \item \textbf{建議}：短途或臨時用車，手機支架夠用；日常通勤長途，建議升級支援CarPlay/HiCar嘅車機或加裝後裝車機，體驗同安全都更好。
\end{itemize}

\begin{tcbwarn}{安全提醒}
無論用邊種方案，揸車時操作手機都違反交規。手機支架淨係可以用嚟「睇導航」，唔可以用嚟「操作手機」。設定目的地、切換歌曲等操作，請務必停車後再做，或使用語音控制。

\end{tcbwarn}

\booksection{六、實用技巧}

\begin{itemize}
  \item \textbf{無線CarPlay連接唔穩定}：(1)關閉車機藍牙再重新配對；(2)刪除車機端已配對設備記錄後重連；(3)檢查手機係咪同時連咗多個藍牙設備造成衝突；(4)升級車機韌體同手機系統到最新版；(5)真係唔得，改用有線CarPlay，更穩定。
  \item \textbf{車機流量用完}：(1)用手機開熱點，車機連手機熱點上網；(2)喺車機設定中將「高耗流量」應用（如影片、線上音樂）切到Wi-Fi；(3)定期清理車機離線地圖快取。
  \item \textbf{駕駛模式自動觸發}：(1)iOS「專注模式」設定「連接車載藍牙時自動開啟」，自動屏蔽非緊急通知；(2)安卓「駕駛模式」或第三方App（如Tasker）可實現連藍牙自動靜音、自動回覆消息；(3)華為手機可設定「智感支付」，上車自動切到車機播放音樂。
\end{itemize}

\begin{tcbinfo}{選車時嘅小建議}
如果喺乎車機體驗，選車時務必去4S店實測：(1)帶自己嘅手機試CarPlay/HiCar連接係咪穩定；(2)試語音助手嘅反應速度同指令識別率；(3)睇應用市場有冇你常用嘅App；(4)關注後續OTA升級頻率——好車機係用軟件迭代出嚟嘅。

\end{tcbinfo}

\begin{tcbpractice}{動手：優化你嘅車機互聯體驗}
\begin{enumerate}
  \item 確認你嘅手機同車機支援嘅互聯方案（CarPlay / HiCar / Android Auto / CarLife）。
  \item 清理一次車機嘅藍牙配對列表，淨係保留常用手機，減少連接衝突。
  \item 設定一個「駕駛專注模式」，連接車機藍牙時自動開啟。
  \item 下載車機品牌App（如鴻蒙智行、理想汽車、蔚來App），睇嚇能否遠程控制冷氣、查看車況。
\end{enumerate}

\end{tcbpractice}

\begin{tcbsummary}{本章小結}
\begin{itemize}
  \item 車機互聯三大路線：投屏（CarPlay等）、深度融合（鴻蒙座艙等）、CarPlay 2.0深度集成
  \item CarPlay生態最廣；HiCar國內體驗好；Android Auto國內需折騰；CarLife適合舊車後裝
  \item 鴻蒙座艙嘅超級桌面、應用接力、共用帳號代表咗深度融合方向
  \item 蔚來NOMI、理想多屏、小米人車家、比亞迪旋轉屏各有特色
  \item 手機當車機僅適合臨時場景，長期建議升級或後裝車機
  \item 實用技巧：CarPlay重連、手機熱點補流量、藍牙觸發駕駛模式
\end{itemize}

\end{tcbsummary}
